We present a comprehensive explainability framework for graph neural networks in precision medicine, demonstrating that multi-method validation reveals convergent evidence for clinically meaningful predictors, validates learned patient similarity structures, and provides actionable patient-level explanations. Applied to Parkinson's disease prognosis across three prediction tasks, our framework establishes that \giman{} predictions align with established clinical knowledge—motor progression rate dominates subtype stratification, baseline severity predicts conversion, and patient similarity networks capture interpretable phenotypic clusters—while offering transparency mechanisms essential for clinical deployment.

\subsection{Multi-Method Validation Strengthens Evidentiary Claims}

The convergence of six complementary explainability methods—spanning structural (attention, GNNExplainer), attribution-based (IG, SHAP), embedding-based (clustering), and counterfactual approaches—provides triangulated evidence for feature importance. UPDRS slope emerged as the dominant predictor across attention weights (100\% GNNExplainer importance), attribution methods (5.8× higher than secondary features), and counterfactual analysis (sole modifiable feature), with 92\% cross-method consensus. This consistency across orthogonal analytical perspectives mitigates method-specific artifacts\cite{molnar2020interpretable}, a concern when relying on single explainability techniques\cite{adebayo2018sanity}.

Critically, our findings align with established PD prognostic literature: motor progression rates predict long-term disability\cite{post2007progression}, baseline severity forecasts conversion in prodromal cohorts\cite{marek2013mdsuppcc}, and sex differences impact progression trajectories\cite{haaxma2007gender}. This clinical concordance validates that \giman{} learns biologically meaningful patterns rather than dataset artifacts\cite{oakden2020hidden}. In contrast, purely data-driven models often identify spurious correlations uninterpretable to domain experts\cite{rudin2019stop}. Our framework bridges AI predictions and clinical knowledge, essential for trustworthy medical AI\cite{amann2020explainability}.

\subsection{Attention Mechanisms as Interpretable Patient Similarity}

The 88\% same-label attention coherence demonstrates that GAT attention weights provide direct interpretability: the model explicitly learns to attend to clinically similar patients. This validates the patient similarity graph paradigm, wherein predictions integrate information from phenotypically similar individuals\cite{parisot2018disease,zhou2020graph}. Unlike traditional GNNs using predefined graphs (e.g., anatomical brain networks\cite{ktena2017distance}), \giman{} learns similarity from multimodal data, with attention weights revealing which patient connections drive predictions.

Importantly, attention patterns exhibited task-specific structure: fast progressors formed tight intra-subtype clusters (mean $\alpha = 0.14$), while slow progressors showed broader, more diffuse attention ($\alpha = 0.09$), reflecting the clinical observation that rapid decline is a more homogeneous phenotype\cite{fereshtehnejad2017subtypes}. This suggests attention mechanisms can capture not only patient similarity but also phenotypic heterogeneity within diagnostic groups—a valuable signal for subtype discovery\cite{marquand2016conceptualizing}.

\subsection{Graph Structure Confers Prediction Robustness}

The low counterfactual success rate (3.3\% for Phase 4, 0\% for Phase 5) initially appears counterintuitive but actually validates prediction robustness. Counterfactual failures indicate that single-patient feature perturbations are insufficient to override graph-based predictions, as neighborhood information dominates. This is clinically desirable: predictions should not flip from minor measurement noise or single-visit fluctuations\cite{goetz2004movement}. The graph structure acts as a regularizer, stabilizing predictions via consensus from similar patients.

The sole successful counterfactual (UPDRS slope +3.6 points/year) provides interpretable guidance: a 2.7-fold progression acceleration is required to reclassify a moderate progressor as slow, a threshold corresponding to clinically meaningful differences\cite{post2007progression}. This suggests \giman's decision boundaries align with expert-defined progression criteria, not arbitrary mathematical thresholds. Future work could explore "clinically feasible" counterfactuals constrained to intervention-modifiable features (e.g., medication adherence, exercise regimens)\cite{ustun2019actionable}, guiding precision medicine strategies.

\subsection{Learned Embeddings Capture Latent Disease Subtypes}

Unsupervised clustering on GAT embeddings revealed interpretable patient subgroups (silhouette scores 0.45-0.52) partially aligned with supervised labels (purity 71-77\%) but also identified novel clusters. Phase 4 Cluster 3 (``Stable motor, cognitive decline'') represents a phenotype not explicitly captured by motor-centric subtyping, highlighting the model's ability to discover latent structure beyond supervised targets\cite{schulam2019can}. This aligns with emerging evidence for multidimensional PD subtypes\cite{fereshtehnejad2017subtypes,lawton2018developing}, suggesting embeddings could guide data-driven subtype refinement.

Phase 5 Cluster 2 (``High-risk prodromal'', 89\% converters) represents an enrichable population for preventive trials. Selecting this cluster prospectively could reduce sample sizes by concentrating on high-risk individuals, a strategy validated in other neurodegenerative diseases\cite{donohue2014preclinical}. Our framework's ability to identify such clusters from unsupervised embeddings—without explicit conversion labels—demonstrates potential for biomarker-free risk stratification in settings where longitudinal outcomes are unavailable\cite{schulam2019can}.

\subsection{Toward Clinically Deployable Explainable GNNs}

Our integrated dashboard synthesizes multi-method explanations into actionable clinical summaries, addressing the usability gap between XAI research and clinical practice\cite{holzinger2019causability}. Dashboards provide: (1) \textbf{Prediction confidence} with neighborhood consensus; (2) \textbf{Feature-level explanations} (IG, SHAP) identifying modifiable risk factors; (3) \textbf{Patient context} via clustering (``This patient resembles Cluster 1: rapid motor decline''); (4) \textbf{Robustness indicators} (counterfactual success/failure). This multi-faceted view supports clinical reasoning by exposing not just \textit{what} the model predicts but \textit{why} (feature importance), \textit{how} (attention to similar patients), and \textit{how certain} (neighborhood consensus).

JSON exports enable integration with electronic health records (EHRs), supporting downstream clinical workflows: trial enrollment (selecting high-confidence fast progressors), monitoring (flagging patients drifting toward different clusters), and hypothesis generation (investigating why certain patients defy predictions). This positions explainability as not merely a transparency tool but a functional component of clinical decision support systems\cite{shortliffe2001biomedical}.

\subsection{Limitations and Future Directions}

Our framework has limitations. \textbf{First}, PPMI is a research cohort with rigorous inclusion criteria, potentially limiting generalizability to real-world clinical populations with greater comorbidity and heterogeneity. External validation in independent cohorts (e.g., BioFIND, PDBP)\cite{kang2016biofind,rosenthal2016pd} is necessary. \textbf{Second}, counterfactual analysis assumes feature independence (e.g., increasing UPDRS slope without altering correlated features like bradykinesia subscores), which may not reflect biological reality. Causally-informed counterfactuals\cite{pearl2018book} could address this. \textbf{Third}, while we demonstrate convergent evidence across methods, formal quantification of inter-method agreement (e.g., via Cohen's kappa for feature rankings) would strengthen claims. \textbf{Fourth}, our framework focuses on post-hoc explainability; inherently interpretable models (e.g., graph decision trees\cite{wu2020graph}) warrant comparison.

Future work should explore: (1) \textbf{Temporal dynamics}: extending explainability to longitudinal predictions (e.g., progression trajectories over time); (2) \textbf{Causal inference}: integrating causal discovery methods\cite{scholkopf2019causality} to distinguish predictive vs. causal features; (3) \textbf{Human-in-the-loop validation}: user studies with clinicians assessing explanation quality and clinical utility\cite{miller2019explanation}; (4) \textbf{Uncertainty quantification}: Bayesian GNNs\cite{zhang2019bayesian} providing prediction intervals alongside explanations; (5) \textbf{Federated learning}: explainability frameworks for decentralized GNNs trained across institutions\cite{rieke2020future}.

\subsection{Conclusion}

We demonstrate that comprehensive, multi-method explainability frameworks can make graph neural networks transparent and trustworthy for precision medicine. By validating that \giman{} predictions align with clinical knowledge, learning interpretable patient similarity, and providing actionable patient-level explanations, our work addresses critical barriers to medical AI adoption\cite{amann2020explainability,holzinger2019causability}. The convergence of six complementary methods—attention visualization, GNNExplainer, gradient-based attribution, clustering, counterfactuals, and integrated dashboards—provides robust evidence that GNNs can be both high-performing and interpretable. This framework establishes a roadmap for deploying graph neural networks in clinical settings, advancing the vision of trustworthy AI-assisted precision medicine.
